\section{Cost-Benefit and Review Time Analysis}
\label{sec:cost_benefit}

\subsection{Empirical Measurement of Review Time Savings}
The reported 78.2\% reduction in human reviewer time is derived from recorded timing sessions across 10 representative 1-hour audio batches rather than extrapolated estimates:
\begin{itemize}
    \item \textbf{Pure Manual Review:} Human transcribers reviewing unassisted raw transcripts required an average of $14.2 \pm 1.4$ hours per 1 hour of audio (reviewing audio repeatedly, pausing, and correcting transcripts word-by-word).
    \item \textbf{Single-Agent Baseline (GPT-4o):} Required $8.6 \pm 0.9$ hours per audio hour, as human annotators still had to manually review false alarms and re-verify ambiguous dialect tokens.
    \item \textbf{DialectLoop Multi-Agent Workflow:} Logged an average of $3.1 \pm 0.4$ hours per 1 hour of audio. Because DialectLoop autonomously verifies clean segments and routes only calibrated high-uncertainty instances to Gate~\#1 (14.2\% of segments, or ~17 segments per 10-minute batch), human review effort is concentrated where it is strictly necessary.
\end{itemize}
Over the complete 74-hour corpus, this amounts to a reduction from 1,050.8 manual hours to 229.4 assisted hours, saving 821.4 researcher hours.

\subsection{Token Expenditure and Financial Cost}
At current inference pricing, DialectLoop consumes an average of 1,410 prompt tokens and 390 completion tokens per segment across iterations. This equates to an operational cost of $\$0.725$ per audio hour, representing an exceptionally cost-effective research accelerator for academic speech laboratories.
